\documentclass[10pt, a4paper]{article}

% Parametri che modificano il file main.tex
% Le uniche parti da cambiare su main.tex sono:
% - tabella versioni
% - testo file

\def\titolo{Analisi dei \\ \ requisiti} % \\ per andare a capo, \ per spazio
\def\titoloHeader{Analisi dei requisiti}
\def\data{}

% \def\listaComponenti{
% Bresolin G.,
% Campese M.,
% Ciriolo I.,
% Dugo A.,
% Feltrin E.,
% Michelon R.,
% Orlandi G.
% }

\usepackage{style}
\usepackage{headerfooter}

\title{\titolo}
\author{SWEetCode}

\begin{document}

% PRIMA PAGINA
\begin{titlepage}
    \thispagestyle{empty}
    \begin{tikzpicture}[remember picture, overlay]
        % TRIANGOLI
        \draw[fill=secondarycolor, secondarycolor] (current page.north west) -- (current page.south west) -- (8.8, -28);
        \draw[fill=primarycolor, primarycolor] (-3, 5) -- (4, -13.6) -- (11, 5);

        % LOGO
        \node [xshift=-5cm, yshift=25cm] (logo) at (current page.south east) {\includegraphics[width=6.5cm]{img/logo.png}};

        % SWEETCODE - DATE
        \node [anchor=north east, align=right, xshift=-1.2cm, yshift=20.5cm, text=black] (sweetcode) at (current page.south east) {\fontsize{32pt}{36pt}\selectfont SWEetCode};
        \draw[line width=4pt, lightcol] ([xshift=-3cm, yshift=-0.37cm]sweetcode.south west) -- ([yshift=-0.37cm]sweetcode.south east);
        \node [anchor=north east, align=right, xshift=-1.2cm, yshift=18.7cm, text=black] (date) at (current page.south east){\fontsize{24pt}{24pt} \selectfont \data};

        % NOME FILE
        \node [anchor=north east, text width=15cm, align=right, xshift=-1.2cm, yshift=17cm, text=black] (titolo) at (current page.south east){\fontsize{48pt}{48pt}\textbf{\titolo}};

        % BOX DATI PARTECIPANTI
        \node[anchor=north east, xshift=-1.2cm, yshift=12cm, minimum width=8cm] (box) at (current page.south east){};

        % COMPONENTI DEL GRUPPO
        \node[anchor=north west, yshift=-1cm, align=left] (dati1) at (box.north west) {\fontsize{16pt}{16pt}\selectfont \textbf{Componenti del gruppo}};
        \draw[line width=4pt, lightcol] (dati1.south west) -- ([xshift=8cm]dati1.south west);
        
        \node[anchor=north west, align=left] (dati11) at (dati1.south west)
        {\fontsize{14pt}{14pt}\selectfont Bresolin G.};
        \node[anchor=north west, yshift=-0.7cm, align=left] (dati12) at (dati1.south west)
        {\fontsize{14pt}{14pt}\selectfont Campese M.};
        \node[anchor=north west, yshift=-1.4cm, align=left] (dati13) at (dati1.south west)
        {\fontsize{14pt}{14pt}\selectfont Ciriolo I.};
        \node[anchor=north west, yshift=-2.1cm, align=left] (dati14) at (dati1.south west)
        {\fontsize{14pt}{14pt}\selectfont Dugo A.};
        \node[anchor=north west, yshift=-2.8cm, align=left] (dati15) at (dati1.south west)
        {\fontsize{14pt}{14pt}\selectfont Feltrin E.};
        \node[anchor=north west, yshift=-3.5cm, align=left] (dati16) at (dati1.south west)
        {\fontsize{14pt}{14pt}\selectfont Michelon R.};
        \node[anchor=north west, yshift=-4.2cm, align=left] (dati17) at (dati1.south west)
        {\fontsize{14pt}{14pt}\selectfont Orlandi G.};
        
        % UNIPD - SWE
        \node [xshift=4.4cm, yshift=2.3cm, draw, secondarycolor, text=white] (uni) at (current page.south west) {\fontsize{20pt}{20pt} \selectfont Università di Padova};
        \node [xshift=0.65cm, yshift=0.7cm, draw, secondarycolor, text=white, below=of uni] (corso) {\fontsize{20pt}{20pt}\selectfont Ingegneria del Software};

        % FIRMA
        % \draw[line width=4pt, lightcol] ([xshift=-1.2cm, yshift=1.8cm]current page.south east) -- ([xshift=-8cm, yshift=1.8cm]current page.south east);
        % \node[anchor=north west, xshift=12.9cm, yshift=1.45cm, align=left] at (current page.south west)
        % {\fontsize{13pt}{13pt}\selectfont L'Amministratore: Feltrin E.};
        
    \end{tikzpicture}
\end{titlepage}

% REGISTRO DELLE VERSIONI
%Registro in ordine dalla più recente alla meno recente!
{\renewcommand{\arraystretch}{1.5}
\section*{Registro delle versioni}

\begin{xltabular}{\textwidth}{c|c|c|c|X}
\label{tab:long}

\textbf{Versione} & \textbf{Data} & \quantities{\textbf{Responsabile di}\\\textbf{stesura}}& \textbf{Revisore} & \quantities{\textbf{Dettaglio e}\\\textbf{motivazioni}} \\
\endfirsthead

\textbf{Versione} & \textbf{Data} & \quantities{\textbf{Responsabile di}\\\textbf{stesura}}& \textbf{Revisore} & \quantities{\textbf{Dettaglio e}\\\textbf{motivazioni}} \\
\endhead

\multicolumn{5}{r}{{Continua nella pagina successiva}} \\
\endfoot

\endlastfoot

\hline
v2.0.0(1) & $2023-11-20$ & \quantities{Bresolin G.\\Campese M.\\Dugo A.\\Michelon R.\\Feltrin E.\\Ciriolo I.} & Orlandi G. & Aggiornata sezione glossario.\\
\hline
v1.0.0(1) & $2023-11-19$ & Michelon R. & Campese M. & Modificata introduzione.\\
\hline
v1.0.1 & $2023-11-14$ & \quantities{Ciriolo I.} & Campese M. &  Studio dei primi casi d'uso.\\
\hline
v0.0.1 & $2023-11-06$ & \quantities{Ciriolo I.\\Orlandi G.} & Campese M. &  Impostazione del documento, Introduzione.\\
    
\end{xltabular}}
\newpage

% INDICE
\tableofcontents
\newpage

% INIZIO PAGINE
\section{Prova Sezione}
\subsection{Prova Sottosezione}

You can't just fit it all in. Also, actions aren't all worth the same. Like in the Pareto principle (80/20), some actions have tremendously better effects than others which are practically useless in the pursuit of your goals. 

\subsubsection{Prova Sottosottosezione}
{\renewcommand{\arraystretch}{1.5}
\begin{tabularx}{\textwidth}{c|X|c|c|c}
\textbf{n°} & \textbf{Azione} & \textbf{Incaricato} & \textbf{Revisore} & \textbf{Scadenza} \\
\hline
5 & Contatto aziende & Ciriolo I. & Team & $2023-10-19$ \\
\hline
6 & Setup workspace informale & Dugo A. & Team & $2023-10-17$ \\
\hline
7 & Creazione template contenuti verbali & Bresolin G. & Team & $2023-10-20$ \\
\hline
8 & Stesura verbale interno 2023-10-17 & Bresolin G. & Feltrin E. & $2023-10-20$ \\
\hline
10 & Creazione logo & Michelon R. & Team & $2023-10-20$ \\
\end{tabularx}}

\section{Seconda Sezione}
\paragraph{Prova paragrafo} Prova testo sotto paragrafo, altro testo per andare a capo e vedere come funziona:
\begin{itemize}
  \item List entries start with the \verb|\item| command.
  \item Individual entries are indicated with a black dot, a so-called bullet.
  \item The text in the entries may be of any length.
\end{itemize}

\end{document}